\documentclass[11pt]{article}

\usepackage[margin=1.02in]{geometry}
\usepackage{amsmath,amssymb,amsthm,mathtools}
\usepackage{booktabs,array}
\usepackage{enumitem}
\usepackage{microtype}
\usepackage[hidelinks]{hyperref}

\allowdisplaybreaks
\numberwithin{equation}{section}
\setlength{\parindent}{1.2em}
\setlength{\parskip}{0.15em}

\newtheorem{theorem}{Theorem}[section]
\newtheorem{lemma}[theorem]{Lemma}
\newtheorem{proposition}[theorem]{Proposition}
\newtheorem{corollary}[theorem]{Corollary}
\theoremstyle{definition}
\newtheorem{remark}[theorem]{Remark}

\newcommand{\per}{\operatorname{per}}
\newcommand{\pcap}{\operatorname{Cap}}
\newcommand{\e}{\mathrm e}
\newcommand{\Rge}{\mathbb R_{\ge0}}
\newcommand{\pos}[1]{#1_{+}}

\title{Face Entropy and Quadratic Tangent Duality\\
for Dittert's Conjecture in Dimension Six}
\author{Standalone dimension-six extension note}
\date{24 July 2026}

\begin{document}
\maketitle

\begin{abstract}
We give a dimension-six extension of the Hall-deficit proof of Dittert's
conjecture.  Two refinements are needed below dimension seven.  First, the
sharp permanent bound for a rectangular zero face is strengthened by retaining
the entropy of a selected zero-block column relative to that face.  This
settles every two-sided Hall configuration of core size at least two after a
critical-core stationarity argument.  Second, in core size one we retain the
exact cofactor entropy, strengthen its quadratic coefficient from $1/2$ to
$5/8$ in dimensions three and four, and eliminate the remaining entropy and
deficit variables with one quadratic supporting line of slope $7/6$.  The
resulting four two-variable and six three-variable inequalities are certified
on rational boxes by exact Bernstein subdivision.  The certificate program
uses only integers and \texttt{fractions.Fraction}; no floating-point assertion
enters the proof.  All labels in this note are prefixed by \texttt{n6:} so that
the note can be merged directly with the dimension-seven and dimensions
$8$--$15$ manuscripts.
\end{abstract}

\section{Statement and interface with the common argument}\label{n6:sec:setup}

For a nonnegative $n\times n$ matrix $A$ whose entries have total sum $n$,
write $r_i$ and $c_j$ for its row and column sums and put
\[
 \Phi(A)=\prod_i r_i+\prod_j c_j-\per(A),
 \qquad
 \gamma_n=\frac{n!}{n^n}.
\]
The result of this note is the following.

\begin{theorem}\label{n6:thm:main}
For every nonnegative $6\times6$ matrix $A$ whose entries sum to $6$,
\[
 \Phi(A)\le 2-\frac{6!}{6^6}=2-\frac5{324}.
\]
Equality holds only at $A=J_6/6$.
\end{theorem}

We use the arbitrary-dimensional infrastructure proved in the common
manuscript.  It is convenient to record precisely the interface.  If $A$ is a
global maximizer, set
\[
 P:=\prod_i r_i=1-\rho,
 \qquad
 Q:=\prod_j c_j=1-\sigma,
 \qquad
 H:=\per(A)=\gamma_n-\delta.
\]
Then
\begin{equation}
 0\le \rho,\sigma,
 \qquad
 \rho+\sigma\le\delta\le\gamma_n.\label{n6:eq:budget}
\end{equation}
In particular,
\begin{equation}
 P,Q\ge1-\gamma_n,
 \qquad
 0\le H\le\gamma_n.\label{n6:eq:basic-deficit}
\end{equation}

For a boundary maximizer, the Hall-deficit construction supplies a critical
pair $I,J\subseteq[n]$.  Put
\[
 R=I^c,
 \quad C=J^c,
 \quad a=|R|,
 \quad b=|C|,
 \quad p=n-a-b\ge1,
\]
and define
\[
 \tau=1-\frac{s_A(I,J)}p,
 \qquad z=1-\tau.
\]
When $\tau>0$ and $z>0$, the dilation $S=A/z$ dominates a doubly stochastic
matrix $B$ satisfying
\begin{equation}
 B[R,C]=0.\label{n6:eq:zero-face}
\end{equation}
The critical cut is saturated:
\begin{equation}
 B[I,J]=S[I,J]=\frac{A[I,J]}z.\label{n6:eq:core-saturation}
\end{equation}
Indeed, both $B[I,J]$ and $S[I,J]$ have total mass $p$, and $B\le S$.
The case $\tau=0$ is already impossible for a boundary maximizer: then $A$ is
doubly stochastic, while its zero entry forces a permanent strictly larger
than $\gamma_n$.

We also use the dimension-free fact that every entrywise positive global
maximizer is $J_n/n$.  Thus Theorem~\ref{n6:thm:main} follows once all critical
Hall configurations in dimension six have been excluded.

For later use, set
\[
 v_k=\frac{k!}{k^k}\quad(k\ge1),
 \qquad v_0=1.
\]
The sharp zero-rectangle constant is
\begin{equation}
 L^{\sharp}_{n,a,b}
 =\frac{v_{n-a}v_{n-b}}{v_{n-a-b}}.\label{n6:eq:sharp-constant}
\end{equation}

\section{Permanent entropy on a zero face}\label{n6:sec:entropy}

The sharp constant \eqref{n6:eq:sharp-constant} remembers only the dimensions
of the zero rectangle.  In dimension six we must also remember how the mass
is distributed along the critical core.

For a probability vector $y=(y_1,\dots,y_m)$, define
\begin{equation}
 \Psi_m(y)
 =\sum_{i=1}^m(1-y_i)\log(1-y_i)
  -m\left(1-\frac1m\right)\log\left(1-\frac1m\right),\label{n6:eq:Psi}
\end{equation}
with the continuous convention at $y_i=1$, and put
\[
 s_m(y)=\sum_{i=1}^m\left(y_i-\frac1m\right)^2.
\]

\subsection{A selected column relative to the zero face}

\begin{theorem}[Selected face-column entropy]\label{n6:thm:face-column}
Let $B$ be an $n\times n$ doubly stochastic matrix with
\[
 B[R,C]=0,
 \qquad |R|=a,
 \qquad |C|=b,
 \qquad p=n-a-b\ge1.
\]
Put
\[
 I=R^c,
 \qquad m=n-a,
 \qquad \ell=n-b.
\]
Choose any column $c\in C$, and let
\[
 y=(b_{ic})_{i\in I}
\]
be that column on its support.  Then
\begin{equation}
 \per(B)
 \ge
 \frac{v_m v_\ell}{v_p}\exp\{\Psi_m(y)\}
 \ge
 \frac{v_m v_\ell}{v_p}
 \exp\left\{\frac12s_m(y)\right\}.\label{n6:eq:face-column}
\end{equation}
\end{theorem}

\begin{proof}
Let
\[
 P_B(x)=\prod_{i=1}^n\left(\sum_{j=1}^n b_{ij}x_j\right).
\]
As in the common capacity argument, $P_B$ is real stable and
$\pcap(P_B)=1$.  Differentiate first in the selected variable $x_c$ and set
$x_c=0$.  If $L_i$ denotes the remaining row form, the derivative is
\[
 q=\sum_{i\in I}y_i\prod_{k\ne i}L_k.
\]
Weighted AM--GM gives
\[
 q\ge\prod_{i=1}^nL_i^{1-y_i},
\]
where $y_i=0$ for $i\in R$.  A second weighted AM--GM inequality in each row
form gives
\[
 \pcap(q)\ge\prod_{i\in I}(1-y_i)^{1-y_i}.\label{n6:eq:q-capacity}
\]

There remain $b-1$ variables indexed by $C\setminus\{c\}$.  Every one is
supported on the $m$ rows in $I$, and the first differentiation has already
removed one of those row factors in every surviving product-rule term.  After
$k$ further zero-block differentiations, the degree of the next such variable
is therefore at most $m-1-k$.  The derivative-capacity inequality supplies
\[
 G(m-1)G(m-2)\cdots G(p+1)=\frac{v_{m-1}}{v_p}.
\]
After all $b$ zero-block variables have been extracted, square-free
coefficient extraction in the remaining $\ell$ variables supplies $v_\ell$.
Consequently,
\[
 \per(B)
 \ge
 \frac{v_{m-1}v_\ell}{v_p}
 \prod_{i\in I}(1-y_i)^{1-y_i}.
\]
Since
\[
 v_{m-1}
 \exp\left\{m\left(1-\frac1m\right)
                   \log\left(1-\frac1m\right)\right\}
 =v_m,
\]
the first inequality in \eqref{n6:eq:face-column} follows.  The second follows
from the $1$-strong convexity of
\[
 h(t)=(1-t)\log(1-t),
 \qquad h''(t)=\frac1{1-t}\ge1.
\]
\end{proof}

\subsection{Exact geometric cofactor entropy}

The codimension-one case permits both core marginals to be retained
simultaneously.

\begin{lemma}[Exact geometric cofactor entropy]\label{n6:lem:cofactor}
Let $M$ be an $m\times m$ doubly stochastic matrix.  Select one column
$\xi=(\xi_1,\dots,\xi_m)$, and let $d_i$ be the permanent of the matrix
obtained by deleting row $i$ and the selected column.  Then
\begin{equation}
 \prod_{i=1}^m d_i^{\xi_i}
 \ge v_m\exp\{\Psi_m(\xi)\}.\label{n6:eq:cofactor-exact}
\end{equation}
\end{lemma}

\begin{proof}
Let $Y$ be the other $m-1$ columns and, initially assuming $\xi_i<1$, normalize
row $i$ of $Y$ by $1-\xi_i$.  Write the resulting row forms as $L_i$ and put
\[
 p_i=\per N_{-i}.
\]
Then
\begin{equation}
 d_i=\left(\prod_{k\ne i}(1-\xi_k)\right)p_i.\label{n6:eq:di-pi}
\end{equation}
For positive $t_i$, consider
\[
 q_t(x)=\sum_i\xi_i t_i\prod_{k\ne i}L_k(x).
\]
This polynomial is stable, because it is the derivative at $z=0$ of
\[
 \prod_i\bigl(L_i(x)+\xi_i t_i z\bigr).
\]
Two weighted AM--GM inequalities give
\[
 \pcap(q_t)\ge\prod_i t_i^{\xi_i}.
\]
Square-free extraction in $m-1$ variables therefore yields
\[
 \sum_i\xi_i t_i p_i
 \ge v_{m-1}\prod_i t_i^{\xi_i}.
\]
Taking $t_i=p_i^{-1}$ and then passing to boundary cases by continuity gives
\[
 \prod_i p_i^{\xi_i}\ge v_{m-1}.
\]
Taking the $\xi$-weighted geometric mean in \eqref{n6:eq:di-pi},
\[
 \prod_i d_i^{\xi_i}
 \ge v_{m-1}\prod_i(1-\xi_i)^{1-\xi_i}
 =v_m\exp\{\Psi_m(\xi)\}.
\]
\end{proof}

\begin{corollary}[Codimension-one core entropy]\label{n6:cor:codim-one}
Let $B$ be doubly stochastic with an $a\times b$ zero block, where
$a+b=n-1$.  Put $I=R^c$, $J=C^c$, and let $\xi,\eta$ be the row and column
marginals of the core $B[I,J]$, whose total mass is one.  Then
\begin{equation}
 \per(B)
 \ge v_{b+1}v_{a+1}
 \exp\{\Psi_{b+1}(\xi)+\Psi_{a+1}(\eta)\}.\label{n6:eq:codim-one-exact}
\end{equation}
\end{corollary}

\begin{proof}
Every perfect matching uses exactly one core edge.  If $d_i$ and $e_j$ are the
two complementary permanental cofactors, then
\[
 \per(B)=\sum_{i,j}x_{ij}d_i e_j.
\]
Weighted AM--GM with weights $x_{ij}$ gives
\[
 \per(B)\ge\prod_i d_i^{\xi_i}\prod_j e_j^{\eta_j}.
\]
Apply Lemma~\ref{n6:lem:cofactor} to the two adjoining doubly stochastic
matrices.
\end{proof}

\subsection{The low-order coefficient \texorpdfstring{$5/8$}{5/8}}

The universal estimate from strong convexity is
\begin{equation}
 \Psi_m(y)\ge\frac12s_m(y).\label{n6:eq:half-entropy}
\end{equation}
For the pair $(a,b)=(2,3)$ in dimension six, the following improvement is
decisive.

\begin{lemma}[Three- and four-dimensional entropy]\label{n6:lem:five-eighths}
For every probability vector $y$ of length $m=3$ or $m=4$,
\begin{equation}
 \Psi_m(y)\ge\frac58s_m(y).\label{n6:eq:five-eighths}
\end{equation}
\end{lemma}

\begin{proof}
For $0\le t\le1$,
\[
 (1-t)\log(1-t)
 =-t+\sum_{k=2}^{\infty}\frac{t^k}{k(k-1)}.
\]
Writing $p_k=\sum_i y_i^k$, we obtain
\begin{equation}
 \Psi_m(y)
 =\sum_{k=2}^{\infty}
   \frac{p_k-m^{1-k}}{k(k-1)}.\label{n6:eq:power-series-Psi}
\end{equation}
Every summand is nonnegative by convexity.

For $m=3$, Schur's inequality and
$q=y_1y_2+y_2y_3+y_3y_1$ give
\[
 p_3-\frac19-\frac34\left(p_2-\frac13\right)
 =3y_1y_2y_3-\frac32q+\frac7{18}
 \ge\frac{1-3q}{18}\ge0.
\]
Keeping only the $k=2,3$ terms in \eqref{n6:eq:power-series-Psi} proves
\[
 \Psi_3\ge\frac12s_3+\frac16\cdot\frac34s_3=\frac58s_3.
\]

Now let $m=4$ and minimize
\[
 F(y)=\Psi_4(y)-\frac58s_4(y)
\]
on the probability simplex.  If one coordinate is zero, the other three form
a probability vector $u$, and the already proved $m=3$ case gives
\[
 F(y)
 \ge
 3h(1/3)-4h(1/4)-\frac5{96}
 =\log\frac{256}{243}-\frac5{96}>0.\label{n6:eq:boundary-gap}
\]
The last inequality is exact: with $x=13/243$, the alternating logarithm
series gives
\[
 \log(1+x)
 >x-\frac{x^2}{2}+\frac{x^3}{3}-\frac{x^4}{4}
 =\frac5{96}+\frac{3635617}{111577100832}.
\]

It remains to consider an interior minimizer.  Put
\[
 \varphi(t)=h(t)-\frac58(t-1/4)^2.
\]
Then
\[
 \varphi''(t)=\frac1{1-t}-\frac54,
\]
which is negative on $(0,1/5)$ and positive on $(1/5,1)$.  The Lagrange
multiplier equation $\varphi'(y_i)=\lambda$ therefore has at most two roots.
The smaller root cannot occur twice at a local minimum, since transferring a
small amount between two such coordinates has negative second variation.
Hence an interior nonuniform minimizer must, after permutation, have the form
\[
 \left(t,\frac{1-t}{3},\frac{1-t}{3},\frac{1-t}{3}\right),
 \qquad 0<t<\frac14.
\]
For this vector, retain the terms $k=3,\dots,7$ in
\eqref{n6:eq:power-series-Psi}.  Exact expansion gives
\[
 F(y)\ge
 \frac{(4t-1)^2}{89579520}
 \left(
 \begin{aligned}
 &133120t^5+255232t^4+387328t^3\\
 &{}+696752t^2+1028168t+2057
 \end{aligned}
 \right)\ge0.
\]
The uniform vector gives equality.  This completes the proof.
\end{proof}

\section{Dimension-six localization and conditional energy}\label{n6:sec:localization}

For the remainder of the proof,
\[
 n=6,
 \qquad \gamma=\gamma_6=\frac5{324}.
\]

\subsection{A uniform lower bound for every row and column sum}

\begin{lemma}\label{n6:lem:row-floor}
Every row sum and every column sum of a hypothetical maximizer satisfies
\begin{equation}
 r_i,c_j>\frac{423}{500}.\label{n6:eq:row-floor}
\end{equation}
\end{lemma}

\begin{proof}
If, for example, $r_1\le423/500$, then AM--GM on the other five row sums gives
The function $t((6-t)/5)^5$ is increasing on $[0,1]$, and hence
\[
 P\le
 \frac{423}{500}
 \left(\frac{6-423/500}{5}\right)^5.
\]
But
\[
 1-\gamma-
 \frac{423}{500}
 \left(\frac{6-423/500}{5}\right)^5
 =
 \frac{29233981969641209}
 {3955078125000000000000}>0,
\]
contradicting $P\ge1-\gamma$.  The column argument is identical.
\end{proof}

\subsection{Direct subset localization}

For $1\le k\le5$, define
\[
 \mathcal P_k(t)
 =\left(1+\frac tk\right)^k
  \left(1-\frac t{6-k}\right)^{6-k}
\]
and set
\begin{equation}
 h_1=\frac{17}{100},
 \quad h_2=\frac{21}{100},
 \quad h_3=\frac{11}{50},
 \quad h_4=\frac15,
 \quad h_5=\frac{31}{200}.\label{n6:eq:h-table}
\end{equation}
The exact gaps are
\begin{equation}
\begin{array}{c|c}
 k&1-\gamma-\mathcal P_k(h_k)\\ \hline
1&\dfrac{99932137192289}{253125000000000000}\\[2mm]
2&\dfrac{38396337601399}{82944000000000000}\\[2mm]
3&\dfrac{7001004061}{11390625000000}\\[2mm]
4&\dfrac{10159}{1296000000}\\[2mm]
5&\dfrac{3512017710706961}{16200000000000000000}
\end{array}\label{n6:eq:subset-gaps}
\end{equation}
All are positive.

\begin{lemma}[Subset boxes]\label{n6:lem:subset-boxes}
If a set of $k$ rows has total sum $k+e$ with $e>0$, then
\begin{equation}
 0<e<h_k.\label{n6:eq:positive-excess}
\end{equation}
For an arbitrary $k$-set,
\begin{equation}
 e>-\frac{77k}{500}.\label{n6:eq:negative-excess}
\end{equation}
The analogous statements hold for columns.
\end{lemma}

\begin{proof}
AM--GM on the set and its complement gives $P\le\mathcal P_k(e)$.
The function $\mathcal P_k$ is strictly decreasing for $e>0$, while
$P\ge1-\gamma$ and \eqref{n6:eq:subset-gaps} holds; this proves
\eqref{n6:eq:positive-excess}.  Equation \eqref{n6:eq:negative-excess}
follows immediately from \eqref{n6:eq:row-floor}.
\end{proof}

For a critical pair, define
\begin{equation}
 e=\sum_{i\in R}r_i-a,
 \qquad
 f=\sum_{j\in C}c_j-b,
 \qquad
 x=s_A(R,C)\ge0.\label{n6:eq:efx}
\end{equation}
Inclusion--exclusion gives the exact relation
\begin{equation}
 p\tau=e+f-x.\label{n6:eq:critical-relation}
\end{equation}
Since $h_k\le11/50$ for every $k$,
\[
 p\tau\le e_++f_+<\frac{11}{25}<1.
\]
Thus $z=1-\tau>0$ in every dimension-six critical configuration, so the Hall
dilation is always available.

\subsection{Conditional inverse-product energy}

For $a\ge1$, put $m=6-a$ and define the groupwise product envelope
\begin{equation}
 P_a(e)
 =\left(1+\frac ea\right)^a
  \left(1-\frac em\right)^m.\label{n6:eq:Pa}
\end{equation}
Thus $P\le P_a(e)$.  If $I=R^c$, define
\[
 V_I=\sum_{i\in I}
 \left(
  \frac1{r_i}-\frac1m\sum_{h\in I}\frac1{r_h}
 \right)^2.
\]

\begin{lemma}[Conditional inverse energy]\label{n6:lem:conditional-energy}
Let
\begin{equation}
 C_*:=\frac{3523}{10000}.\label{n6:eq:Cstar}
\end{equation}
Then
\begin{equation}
 V_I\le\frac{P_a(e)-P}{C_*}.\label{n6:eq:conditional-energy}
\end{equation}
The analogous column inequality holds with $P_b(f)-Q$.
\end{lemma}

\begin{proof}
Put $q=1-e/m$, the average of the row sums in $I$.  By
Lemma~\ref{n6:lem:row-floor}, both $r_i$ and $q$ exceed $423/500$.  For any
$r,q\ge423/500$,
\begin{equation}
 \frac rq-1-\log\frac rq
 \ge
 \frac{(423/500)^2}{2}
 \left(\frac1r-\frac1q\right)^2.\label{n6:eq:scalar-energy}
\end{equation}
To see this, put $y=1-q/r$.  The left side is
\[
 \int_0^y\frac{t}{(1-t)^2}\,dt.
\]
For $y\ge0$, compare the integrand with $t$.  For $y<0$, compare on
$[y,0]$ with the constant denominator $(1-y)^2$ and use
$r=q/(1-y)\ge423/500$.

The linear terms cancel inside each fixed-sum group, so
\[
 \log\frac{P_a(e)}P
 \ge\frac{(423/500)^2}{2}V_I.
\]
Writing $L=\log(P_a/P)\ge0$,
\[
 P_a-P=P(\e^L-1)\ge P L\ge(1-\gamma)L.
\]
Finally,
\[
 \frac{(1-\gamma)(423/500)^2}{2}-C_*
 =\frac{71}{2000000}>0.
\]
This proves \eqref{n6:eq:conditional-energy}.
\end{proof}

The same proof with the full set of rows gives the global energy estimate
\begin{equation}
 \rho\ge C_*\sum_{i=1}^6\left(1-\frac1{r_i}\right)^2,\label{n6:eq:global-energy}
\end{equation}
and similarly for columns.  The inequality is strict whenever the row sums
are not all equal; this is the only regime in which it is used below.

\section{Critical-core stationarity}\label{n6:sec:stationarity}

Let $\xi=(\xi_i)_{i\in I}$ and $\eta=(\eta_j)_{j\in J}$ be the row and column
marginals of the saturated core $B[I,J]$.  Thus
\[
 \sum_{i\in I}\xi_i=\sum_{j\in J}\eta_j=p.
\]
Put
\begin{equation}
 s_I=\sum_{i\in I}\left(\xi_i-\frac p{6-a}\right)^2,
 \qquad
 s_J=\sum_{j\in J}\left(\eta_j-\frac p{6-b}\right)^2,
 \qquad s=s_I+s_J.\label{n6:eq:core-s}
\end{equation}

Let $H_k$ denote the total permanent weight of permutations using exactly
$k$ entries of the forbidden rectangle $R\times C$.  Then
\[
 H=\sum_{k=0}^{d}H_k,
 \qquad d=\min(a,b).
\]
A permutation using $k$ forbidden edges uses exactly $p+k$ core edges.

\begin{lemma}[Critical-core stationarity identity]\label{n6:lem:core-stationarity}
At a global maximizer,
\begin{equation}
 P\left(\sum_{i\in I}\frac{\xi_i}{r_i}-p\right)
 +Q\left(\sum_{j\in J}\frac{\eta_j}{c_j}-p\right)
 =\frac{p\tau H+\sum_{k=1}^{d}kH_k}{z}.\label{n6:eq:core-stationarity}
\end{equation}
\end{lemma}

\begin{proof}
Multiply the saturated core $A[I,J]$ by a parameter $t$, and then multiply the
whole matrix by
\[
 \lambda(t)=\frac6{6+(t-1)pz}
\]
to keep the total entry sum equal to six.  Before global normalization, the
derivatives of the two product terms are
\[
 zP\sum_{i\in I}\frac{\xi_i}{r_i},
 \qquad
 zQ\sum_{j\in J}\frac{\eta_j}{c_j},
\]
while the permanent derivative is
\[
 pH+\sum_{k=1}^{d}kH_k.
\]
The derivative of $\lambda(t)^6$ at $t=1$ is $-pz$.  Stationarity and
$\Phi=P+Q-H$ now give \eqref{n6:eq:core-stationarity}.
\end{proof}

Let $H_0$ denote the contribution of permutations avoiding $R\times C$.
Since $B\le A/z$ and $B[R,C]=0$,
\begin{equation}
 H_0\ge z^6\per(B).\label{n6:eq:H0-B}
\end{equation}
Also,
\begin{equation}
 \sum_{k=1}^{d}kH_k\le d(H-H_0).\label{n6:eq:Hk-bound}
\end{equation}

For the row and column product envelopes, put
\begin{equation}
 u_r=P_a(e)-P,
 \qquad
 u_c=P_b(f)-Q,\label{n6:eq:uruc}
\end{equation}
and define
\begin{equation}
 D=2-P_a(e)-P_b(f),
 \qquad
 K=\gamma-D,
 \qquad
 u=\delta-D.\label{n6:eq:DKu}
\end{equation}
The joint deficit budget gives
\begin{equation}
 u\ge u_r+u_c\ge0,
 \qquad
 H=K-u.\label{n6:eq:u-budget}
\end{equation}

Write
\[
 m=6-a,
 \qquad \ell=6-b,
 \qquad
 E=\frac e{m-e},
 \qquad
 F=\frac f{\ell-f}.
\]
The harmonic-mean inequality and Lemma~\ref{n6:lem:conditional-energy} imply
\begin{align}
 P\left(\sum_{i\in I}\frac{\xi_i}{r_i}-p\right)
 &\ge
 P_a(e)\,pE-u_rpE-\sqrt{\frac{s_Iu_r}{C_*}},\label{n6:eq:row-core-lower}\\
 Q\left(\sum_{j\in J}\frac{\eta_j}{c_j}-p\right)
 &\ge
 P_b(f)\,pF-u_cpF-\sqrt{\frac{s_Ju_c}{C_*}}.\label{n6:eq:column-core-lower}
\end{align}
Indeed, subtract the mean of $(1/r_i)_{i\in I}$ from the zero-sum vector
$(\xi_i-p/m)_{i\in I}$ and use Cauchy--Schwarz.  Adding the two inequalities,
\begin{equation}
 \text{left side of \eqref{n6:eq:core-stationarity}}
 \ge
 P_a(e)pE+P_b(f)pF-u_rpE-u_cpF
 -\sqrt{\frac{su}{C_*}}.\label{n6:eq:core-lower-combined}
\end{equation}

\section{Two-sided faces with \texorpdfstring{$p\ge2$}{p >= 2}}\label{n6:sec:pge2}

Up to transposition, the possible pairs are
\begin{equation}
 (a,b)=(1,1),(1,2),(1,3),(2,2).\label{n6:eq:pge2-pairs}
\end{equation}
Put
\[
 m=6-a,
 \qquad \ell=6-b,
 \qquad p=6-a-b,
 \qquad L=L^{\sharp}_{6,a,b}.
\]

\subsection{From face-column entropy to core entropy}

The $b$ columns in $C$ are probability vectors on $I$.  Their average is
\[
 \left(\frac{1-\xi_i}{b}\right)_{i\in I}.
\]
Since $m=p+b$,
\[
 \frac{1-\xi_i}{b}-\frac1m
 =-\frac1b\left(\xi_i-\frac pm\right).
\]
Convexity of squared distance therefore shows that one of those columns has
squared deviation at least $s_I/b^2$.  Theorem~\ref{n6:thm:face-column} gives
\[
 \per(B)\ge L\exp\left(\frac{s_I}{2b^2}\right).
\]
After transposition,
\[
 \per(B)\ge L\exp\left(\frac{s_J}{2a^2}\right).
\]
Consequently,
\begin{equation}
 \per(B)
 \ge L\exp\left(\frac{s}{2(a^2+b^2)}\right).\label{n6:eq:pge2-core-entropy}
\end{equation}

Put
\begin{equation}
 t=e+f,
 \qquad
 z_0=1-\frac tp,
 \qquad
 A_0=Lz_0^6.\label{n6:eq:pge2-A0}
\end{equation}
Criticality gives $p\tau=t-x\le t$, hence $z\ge z_0$.  Equations
\eqref{n6:eq:H0-B} and \eqref{n6:eq:pge2-core-entropy} yield
\begin{equation}
 H_0\ge A_0
 \exp\left(\frac{s}{2(a^2+b^2)}\right).\label{n6:eq:pge2-H0}
\end{equation}

By \eqref{n6:eq:Hk-bound}, the right side of the stationarity identity is at
most
\[
 \frac{(t+d)H-dH_0}{z_0}.
\]
Define
\begin{equation}
 \mathcal B
 =P_a(e)pE+P_b(f)pF
 -\frac{(t+d)K-dA_0}{z_0}.\label{n6:eq:pge2-B}
\end{equation}
On the localized boxes, $pE_+,pF_+<1$, while criticality gives
$(t+d)/z_0>d\ge1$.  Therefore the terms $u_rpE+u_cpF$ in
\eqref{n6:eq:core-lower-combined} are absorbed by the negative
$-(t+d)u/z_0$ term from the stationarity upper bound.  The remaining entropy
term is also nonpositive.  Thus a necessary condition is
\begin{equation}
 \mathcal B\le\sqrt{\frac{su}{C_*}}.\label{n6:eq:pge2-B-energy}
\end{equation}

At the same time, \eqref{n6:eq:pge2-H0} and $H=K-u$ give
\[
 u+A_0\exp\left(\frac{s}{2(a^2+b^2)}\right)\le K.
\]
If $A_0>K$, this is already impossible.  Otherwise,
$\e^y\ge1+y$ and AM--GM imply
\begin{equation}
 su\le
 \frac{(a^2+b^2)(K-A_0)^2}{2A_0}.\label{n6:eq:pge2-su-upper}
\end{equation}
Consequently, whenever $\mathcal B>0$, stationarity is impossible if
\begin{equation}
 2A_0C_*\mathcal B^2
 >(a^2+b^2)(K-A_0)^2.\label{n6:eq:pge2-master}
\end{equation}

\subsection{Exact bivariate certificate}

Lemma~\ref{n6:lem:subset-boxes} places $(e,f)$ in the rational rectangle
\begin{equation}
 -\frac{77a}{500}\le e\le h_a,
 \qquad
 -\frac{77b}{500}\le f\le h_b.\label{n6:eq:pge2-box}
\end{equation}
This rectangle is deliberately larger than the actual critical region.

Let
\[
 Q_0=(m-e)(\ell-f)z_0,
 \qquad
 \widetilde{\mathcal B}=Q_0\mathcal B,
\]
and clear the denominator in \eqref{n6:eq:pge2-master}:
\begin{equation}
 \mathcal C
 =2A_0C_*\widetilde{\mathcal B}^{\,2}
 -(a^2+b^2)(K-A_0)^2Q_0^2.\label{n6:eq:pge2-C}
\end{equation}
All three functions $A_0-K$, $\widetilde{\mathcal B}$, and $\mathcal C$ are
polynomials with rational coefficients.

Exact Bernstein subdivision proves on the whole rectangle
\eqref{n6:eq:pge2-box} that every point satisfies one of the alternatives
\[
 A_0-K>0,
 \qquad\text{or}\qquad
 \widetilde{\mathcal B}>0\ \text{ and }\ \mathcal C>0.
\]
The first is the direct permanent contradiction; the second is
\eqref{n6:eq:pge2-master}.  The subdivision audit is
\begin{equation}
\begin{array}{c|c|c|c|c}
(a,b)&\text{tree nodes}&\text{maximum depth}&
 A_0-K\text{ leaves}&\text{stationarity leaves}\\ \hline
(1,1)&65&12&26&7\\
(1,2)&51&12&19&7\\
(1,3)&89&13&30&15\\
(2,2)&83&12&27&15
\end{array}\label{n6:eq:pge2-audit}
\end{equation}
Every coefficient transformation is exact.  The program first computes
rational Bernstein coefficients, clears each tensor to integers, and then
uses integer de Casteljau midpoint subdivision.  Hence this is a continuum
certificate, not a grid computation.  All two-sided faces with $p\ge2$ are
excluded.

\section{Codimension-one faces: \texorpdfstring{$p=1$}{p = 1}}\label{n6:sec:p1}

Up to transposition, only
\begin{equation}
 (a,b)=(1,4),(2,3)\label{n6:eq:p1-pairs}
\end{equation}
remain.  Here
\[
 \tau=e+f-x,
 \qquad
 z=1-e-f+x.
\]
The core has total mass one.  Put $s=s_I+s_J$ as before and define
\begin{equation}
 \kappa_{1,4}=\frac12,
 \qquad
 \kappa_{2,3}=\frac58.\label{n6:eq:kappa-pairs}
\end{equation}
Corollary~\ref{n6:cor:codim-one},
\eqref{n6:eq:half-entropy}, and
Lemma~\ref{n6:lem:five-eighths} give
\begin{equation}
 H_0\ge A_0\e^{\kappa s},
 \qquad
 A_0=v_{6-a}v_{6-b}z^6,\label{n6:eq:p1-H0}
\end{equation}
where $\kappa$ is the value corresponding to the pair.

\subsection{The six rational boxes}

Criticality gives $e+f>0$ and $0\le x<e+f$.  The positive-excess bounds,
Lemma~\ref{n6:lem:row-floor}, and $e+f>0$ place every case in one of the
following boxes.  The interval for $x$ is enlarged to a rectangular upper
bound.
\begin{equation}
\begin{array}{c|c|c|c|c}
(a,b)&\text{signs}&e\text{ interval}&f\text{ interval}&x\text{ interval}\\ \hline
(1,4)&++&[0,17/100]&[0,1/5]&[0,37/100]\\
(1,4)&-+&[-77/500,0]&[0,1/5]&[0,1/5]\\
(1,4)&+-&[0,17/100]&[-17/100,0]&[0,17/100]\\
(2,3)&++&[0,21/100]&[0,11/50]&[0,43/100]\\
(2,3)&-+&[-11/50,0]&[0,11/50]&[0,11/50]\\
(2,3)&+-&[0,21/100]&[-21/100,0]&[0,21/100]
\end{array}\label{n6:eq:p1-boxes}
\end{equation}
Throughout all six boxes, $z>0$, $m-e>0$, and $\ell-f>0$.

\subsection{Stationarity with the entropy and deficit variables retained}

Keep the definitions $D,K,u$ from \eqref{n6:eq:DKu} and put
\[
 E=\frac e{m-e},
 \qquad
 F=\frac f{\ell-f},
 \qquad
 d=\min(a,b).
\]
Comparison of \eqref{n6:eq:core-lower-combined} with
\eqref{n6:eq:core-stationarity}, \eqref{n6:eq:Hk-bound}, and
\eqref{n6:eq:p1-H0} gives
\begin{equation}
 \mathcal B
 \le
 \sqrt{\frac{su}{C_*}}
 +Wu
 -\frac{dA_0}{z}(\e^{\kappa s}-1),\label{n6:eq:p1-stationarity}
\end{equation}
where
\begin{align}
 \mathcal B
 &=P_a(e)E+P_b(f)F
 -\frac{(\tau+d)K-dA_0}{z},\label{n6:eq:p1-B}\\
 W&=E_++F_+-\frac{\tau+d}{z}.\label{n6:eq:p1-W}
\end{align}
Here $E_+$ and $F_+$ are chosen according to the sign branch in
\eqref{n6:eq:p1-boxes}.  The estimate uses
$u_rE+u_cF\le u(E_++F_+)$.

The permanent budget is
\begin{equation}
 u+A_0\e^{\kappa s}\le K.\label{n6:eq:p1-per-budget}
\end{equation}
Set
\begin{equation}
 v=A_0\kappa s,
 \qquad
 q=\frac dz.\label{n6:eq:vq}
\end{equation}
The second-order exponential bound
$\e^y\ge1+y+y^2/2$ transforms
\eqref{n6:eq:p1-stationarity} and \eqref{n6:eq:p1-per-budget} into
\begin{align}
 u+v+\frac{v^2}{2A_0}&\le R_0:=K-A_0,\label{n6:eq:p1-budget-v}\\
 \mathcal B&\le
 \sqrt{\frac{uv}{A_0\kappa C_*}}
 +Wu-q\left(v+\frac{v^2}{2A_0}\right).\label{n6:eq:p1-stationarity-v}
\end{align}

\subsection{One quadratic tangent and an exact cone maximum}

Choose
\begin{equation}
 \lambda=\frac16,
 \qquad c=1+\lambda=\frac76.\label{n6:eq:tangent-lambda}
\end{equation}
For every $v\ge0$,
\begin{equation}
 v+\frac{v^2}{2A_0}
 \ge cv-\frac{A_0\lambda^2}{2},\label{n6:eq:quadratic-tangent}
\end{equation}
because the difference is $(v-A_0\lambda)^2/(2A_0)$.  Define
\begin{equation}
 R_\lambda=K-A_0+\frac{A_0\lambda^2}{2},
 \qquad
 \mathcal B_\lambda
 =\mathcal B-\frac{qA_0\lambda^2}{2},\label{n6:eq:Rlambda-Blambda}
\end{equation}
and put $w=cv$.  Equations
\eqref{n6:eq:p1-budget-v}--\eqref{n6:eq:quadratic-tangent} imply the necessary
condition
\begin{equation}
 \mathcal B_\lambda
 \le
 \sqrt{\frac{uw}{cA_0\kappa C_*}}+Wu-qw,
 \qquad
 u,w\ge0,
 \quad u+w\le R_\lambda.\label{n6:eq:triangle-problem}
\end{equation}

The exact maximum of the right side over this triangle is
\begin{equation}
 \max\left\{
 0,
 \frac{R_\lambda}{2}
 \left[
  W-q+\sqrt{(W+q)^2+\frac1{cA_0\kappa C_*}}
 \right]
 \right\}.\label{n6:eq:triangle-maximum}
\end{equation}
Indeed, homogeneity places every positive maximum on $u+w=R_\lambda$;
writing $u=R_\lambda t$, $w=R_\lambda(1-t)$ reduces the problem to the largest
eigenvalue of a real symmetric $2\times2$ matrix.

Define
\begin{equation}
 F_1=\mathcal B_\lambda+qR_\lambda,
 \qquad
 F_2=\mathcal B_\lambda-WR_\lambda.\label{n6:eq:F1F2}
\end{equation}
The following four strict inequalities contradict
\eqref{n6:eq:triangle-maximum}:
\begin{equation}
 \mathcal B_\lambda>0,
 \qquad F_1>0,
 \qquad F_2>0,
 \qquad
 4cA_0\kappa C_*F_1F_2>R_\lambda^2.\label{n6:eq:tangent-criterion}
\end{equation}
The first inequality beats the zero vertex.  The middle two make the quantity
that is squared positive, and the final inequality is exactly the squared
comparison with the second term in \eqref{n6:eq:triangle-maximum}.

\subsection{Exact trivariate certificate}

Let
\[
 M=(m-e)(\ell-f),
 \qquad
 Q_0=Mz.
\]
Clear the positive denominator $Q_0$ in
$\mathcal B_\lambda$, $F_1$, $F_2$, and clear $Q_0^2$ in the last inequality
of \eqref{n6:eq:tangent-criterion}.  The resulting four expressions are
polynomials in $(e,f,x)$ with rational coefficients.

Exact Bernstein subdivision proves on each full box in
\eqref{n6:eq:p1-boxes} that every point satisfies either
\[
 A_0-K>0,
\]
or all four cleared inequalities in \eqref{n6:eq:tangent-criterion}.  Thus the
certificate is stronger than necessary: it does not need to discard the
noncritical part $x\ge e+f$.  Its audit is
\begin{equation}
\begin{array}{c|c|c|c|c|c}
(a,b)&\text{signs}&\text{tree nodes}&\text{maximum depth}&
 A_0-K\text{ leaves}&\text{tangent leaves}\\ \hline
(1,4)&++&485&18&130&113\\
(1,4)&-+&59&13&23&7\\
(1,4)&+-&89&15&38&7\\
(2,3)&++&427&17&116&98\\
(2,3)&-+&51&14&21&5\\
(2,3)&+-&75&14&31&7
\end{array}\label{n6:eq:p1-audit}
\end{equation}
As in Section~\ref{n6:sec:pge2}, all sign decisions are made on
integer-cleared Bernstein tensors.  This excludes both codimension-one Hall
faces.

\section{One-sided Hall cuts}\label{n6:sec:one-sided}

It remains to exclude a critical pair with, after transposition, $a=0$.  Then
$I=[6]$, $|C|=b$, $p=6-b$, and \eqref{n6:eq:critical-relation} becomes
\[
 p\tau=f.
\]
Lemma~\ref{n6:lem:subset-boxes} gives
\begin{equation}
 0<\tau<\max_{1\le b\le5}\frac{h_b}{6-b}
 =\frac{31}{200}.\label{n6:eq:one-sided-tau}
\end{equation}

The common one-sided stationarity argument applies verbatim.  If $\beta$ is
the average row distribution of the saturated $p$ columns and
\[
 s=\sum_{i=1}^6\left(\beta_i-\frac16\right)^2,
\]
then
\begin{equation}
 D:=1-\sum_i\frac{\beta_i}{r_i}
 \ge\frac{(1-\gamma)\tau}{1-\tau}>0.\label{n6:eq:one-sided-D}
\end{equation}
Combining Cauchy--Schwarz with the global energy
\eqref{n6:eq:global-energy},
\begin{equation}
 \rho>
 C_*\frac{(1-\gamma)^2\tau^2}{(1-\tau)^2s}.\label{n6:eq:one-sided-rho}
\end{equation}
At least one of the saturated columns has squared deviation at least $s$.
The column-entropy permanent bound and the inherited zero of $B$ give
\begin{equation}
 \per(B)\ge\max\{\kappa_6,\gamma\e^{s/2}\},
 \qquad
 \kappa_6=v_5G(5)=\frac{6144}{390625}>\gamma.\label{n6:eq:one-sided-perB}
\end{equation}
Hence
\begin{equation}
 H\ge(1-\tau)^6\max\{\kappa_6,\gamma\e^{s/2}\},
 \qquad
 H\le\gamma-\rho.\label{n6:eq:one-sided-sandwich}
\end{equation}

Put $s_0=2\log(\kappa_6/\gamma)$.  When $s\le s_0$, the same completion of the
square as in the common proof reduces the contradiction to
\begin{equation}
 C_*(1-\gamma)^2\gamma-\frac{6^2}{2}\kappa_6^2
 =
 \frac{67858438997036543}
 {83037656250000000000}>0.\label{n6:eq:small-entropy}
\end{equation}

When $s\ge s_0$, use $\e^{s/2}\ge1+s/2$ and optimize the two terms in $s$.
With $t=\tau$ and $z=1-t$, it is enough to prove
\begin{equation}
 2C_*(1-\gamma)^2(1-t)^4
 -\gamma\left(6-15t+20t^2-15t^3+6t^4-t^5\right)^2>0\label{n6:eq:large-poly}
\end{equation}
on $0\le t\le31/200$.  Every Bernstein coefficient of this degree-ten
polynomial on that interval is positive.  The smallest coefficient is
\begin{equation}
 \frac{47529302914242746885557519}
 {537477120000000000000000000}>0.\label{n6:eq:large-min}
\end{equation}
Thus all one-sided Hall cuts are impossible.

\section{Completion}\label{n6:sec:completion}

Let $A$ be a hypothetical boundary global maximizer in dimension six.  The
case $\tau=0$ is excluded by the common Hall argument.  When $\tau>0$, the
localization in Section~\ref{n6:sec:localization} makes the Hall dilation
valid.  One-sided cuts are excluded by Section~\ref{n6:sec:one-sided}.
Two-sided cuts with $p\ge2$ are excluded by the four exact certificates in
Section~\ref{n6:sec:pge2}, and the two codimension-one cuts are excluded by the
six certificates in Section~\ref{n6:sec:p1}.  Hence no boundary global
maximizer exists.

Every global maximizer is therefore positive, and the dimension-free
positive-maximizer theorem gives $A=J_6/6$.  This proves
Theorem~\ref{n6:thm:main}.

\begin{remark}[Merge point]\label{n6:rem:merge}
The only ingredients specific to dimension six are the constants in
Section~\ref{n6:sec:localization}, the selected face-column entropy used for
$p\ge2$, the $5/8$ low-order cofactor entropy, and the tangent
$\lambda=1/6$.  The deficit budget, Hall dilation, critical-core variation,
and positive-maximizer theorem are common with the dimension-seven and
$8$--$15$ arguments.  Thus a combined manuscript can retain Sections
\ref{n6:sec:entropy}, \ref{n6:sec:localization},
\ref{n6:sec:pge2}, and \ref{n6:sec:p1} as the dimension-six module and move
the common material to a single preliminary part.
\end{remark}

\end{document}
